\documentclass[a4paper, 10pt]{article}
\usepackage[T1, T2A]{fontenc}
\usepackage[utf8]{inputenc} 
\usepackage[english, russian]{babel}
\usepackage{amsfonts,amssymb,amsthm,mathtools,amsmath}
\usepackage{indentfirst} 
\usepackage{graphicx}
\usepackage{hyperref}
\usepackage{biblatex}
\usepackage{geometry}
\addbibresource{lit.bib}


\geometry{
	a4paper,
	total={170mm,257mm},
	left=20mm,
	top=20mm,
}
\theoremstyle{plain}
\newtheorem*{pro}{Доказательство}
\newtheorem*{ex}{Пример}
\newtheorem{lem}{Лемма}
\newtheorem{theorem}{Теорема}
\newtheorem{defin}{Определение}
\newtheorem*{rem}{Замечание}
\newtheorem*{sled}{Следствие}
\newcommand{\defeq}{:=}
\newcommand{\tang}[1][\hspace{0pt}]{\frac{\partial }{\partial #1}}
\newcommand\Lie{Lie}
\author{я}
\title{,,,}
\begin{document}
	\maketitle
	Пусть нам даны весовые соотношения
	\begin{align*}
		\operatorname{wt}(\alpha_{\mu - 1}) = &\dfrac{1}{\mu + 1}\\
		\operatorname{wt}(\alpha_{\mu - 2}) = &\dfrac{2}{\mu + 1} = \operatorname{wt}(s_{\mu - 2})\\
		&\dots\\
		\operatorname{wt}(\alpha_{0}) = & \dfrac{\mu}{\mu + 1} = \operatorname{wt}(s_{0})
	\end{align*}
	\begin{equation}
		e_1^{\circ(N)}
		=
		\sum_{k=0}^{N-1}\alpha_k(s)e_k \quad \text{где } \quad \alpha_k(0)=0.
		\label{eq:fundamental_relation} 
	\end{equation}
	\begin{equation}
		\left.\frac{\partial\alpha_k}{\partial s_r}\right|_0=0 \qquad\text{если }k\ne r.
		\label{eq:weight_diagonal}
	\end{equation}
	\begin{equation}
		e_a\circ e_b=e_{a+b}
		\qquad\text{для }a+b\le N-1.
		\label{eq:low_product}
	\end{equation}
	\begin{theorem}[Vanishing coefficients]
		Пусть \(M\) - неприводимое полупростое в общем многообразие с градуировкой типа \(A_\mu\), для которого выполнено
		\[
		T_pM \simeq \mathbb{C}\left[x, s_0, \dots, s_{\mu - 1}\right]/\left(x^{\mu} - \sum_{k = 0}^{\mu - 2} \pi_k s_{k} \cdot x^k\right), \text{где } \pi_k \in \mathbb{C}
		\]  Тогда \(M\) изоморфно многообразию простой особенности \(A_{\mu}\).
	\end{theorem}
	\begin{proof}
	\begin{equation}
		\begin{aligned}
			&\text{Определим }B_r:=
			\left.\frac{\partial\alpha_r}{\partial s_r}\right|_0,
			\qquad r=1,\ldots,N-2.
			\label{eq:Br_def}
		\end{aligned}
	\end{equation}
	\end{proof}
\end{document}